\ \ \\[2cm]
The growing integration of pets as family members, coupled with the busy routines of owners, drives the search for technological solutions that ensure animal welfare. This work presents the design, validation, and evolution of a low-cost Internet of Things (IoT) based system for monitoring and automating pet environments. The prototype, validated in high-fidelity simulation (TCC 1), uses an ESP32 microcontroller to collect data from multiple sensors and control actuators.

As an evolution (TCC 2), and in response to a Software Engineering requirements analysis, the work implements two advanced extensions. First, a predictive health module, which applies \textit{Machine Learning} techniques (Isolation Forest) on a controlled \textit{dataset} to detect anomalies in consumption patterns, thereby validating the system's intelligence. Second, the architecture is migrated to a scalable and secure cloud platform (AWS IoT Core and Lambda), enabling robust integration and control via voice commands through an Amazon Alexa \textit{Skill}.

Logic orchestration, data visualization on a \textit{dashboard}, and the sending of notifications are managed by the visual programming tool Node-RED. The objective is to provide an integrated solution that ensures the animal's feeding, hydration, and thermal comfort, allowing for remote and voice control, and providing predictive \textit{insights} into the pet's well-being.

\vspace{1cm}
Keywords: Internet of Things; Software Engineering; Anomaly Detection; AWS; Alexa; ESP32; Node-RED; Animal Welfare.